\documentclass[reqno]{amsart}
\usepackage[T1]{fontenc}
\usepackage{amsmath,amssymb,amsthm}
\usepackage{xr-hyper}
\usepackage[hidelinks]{hyperref}
\externaldocument[dt-]{20261125T002000--definite-theta-lifts__workbook}
\externaldocument[tn-]{20261125T002100--definite-theta-lifts-and-norms__workbook}
\externaldocument[da-]{20261126T030300--the-data__workbook}
\externaldocument[bp-]{20261126T020210--bessel-poincare-series-and-their-gram-operators__workbook}

% Chapter number; \ref{chap} gives it.
\newcounter{chap}
% Chapter numbers, generated from 20261125T012000--grh-workbook__guide.org
% by scripts/make-chapters.py.  Do not edit: change the order of the
% index and rerun the script (./preview-site.sh and CI do this).
%
% Usage in a sheet, after \newcounter{chap}:
%   % Chapter numbers, generated from 20261125T012000--grh-workbook__guide.org
% by scripts/make-chapters.py.  Do not edit: change the order of the
% index and rerun the script (./preview-site.sh and CI do this).
%
% Usage in a sheet, after \newcounter{chap}:
%   % Chapter numbers, generated from 20261125T012000--grh-workbook__guide.org
% by scripts/make-chapters.py.  Do not edit: change the order of the
% index and rerun the script (./preview-site.sh and CI do this).
%
% Usage in a sheet, after \newcounter{chap}:
%   \input{chapters}\chapnumber{<sheet ID>}
\makeatletter
\providecommand\chapters@set[2]{\@namedef{chapters@#1}{#2}}
\chapters@set{20261125T001000}{1}% zeros-to-ramp-sums
\chapters@set{20261125T001100}{2}% inert-primes-and-estimate-b
\chapters@set{20261125T005000}{3}% galois-fourier-and-averaging
\chapters@set{20261125T005500}{4}% columns-from-a-sphere
\chapters@set{20261125T000000}{5}% slope-method
\chapters@set{20261126T030100}{6}% a-twist-for-each-prime
\chapters@set{20261125T003000}{7}% the-chain-and-the-square-root
\chapters@set{20261125T010000}{8}% rank-lemma-and-zero-estimate
\chapters@set{20261125T010500}{9}% poincare-series-and-near-orthogonality
\chapters@set{20261124T223000}{10}% elliptic-curves-for-the-proof
\chapters@set{20261125T009500}{11}% fock-model-and-theta-intertwiners
\chapters@set{20261124T223100}{12}% complex-multiplication-and-inert-reduction
\chapters@set{20261125T008000}{13}% one-cm-determinant
\chapters@set{20261124T214500}{14}% supersingular-theta-divisibility
\chapters@set{20261124T214600}{15}% the-isotropic-plane
\chapters@set{20261125T013700}{16}% pairings-on-a-formal-group-of-height-two
\chapters@set{20261125T013500}{17}% hecke-nondegeneracy-on-the-tree
\chapters@set{20261126T030200}{18}% the-bound-is-attained
\chapters@set{20261125T009000}{19}% kronecker-limit-and-chowla-selberg
\chapters@set{20261125T002000}{20}% definite-theta-lifts
\chapters@set{20261125T015500}{21}% siegel-weil-in-the-smallest-case
\chapters@set{20261125T002100}{22}% definite-theta-lifts-and-norms
\chapters@set{20261125T016500}{23}% hecke-operators-and-theta-lifts
\chapters@set{20261125T016700}{24}% divisibility-along-the-chain
\chapters@set{20261125T004000}{25}% denominators-and-sha
\chapters@set{20261125T004100}{26}% tate-twists-and-sha
\chapters@set{20261126T020430}{27}% the-minus-part-by-gauss-sums
\chapters@set{20261125T014100}{28}% congruence-modules
\chapters@set{20261125T014500}{29}% eisenstein-congruences-and-selmer-classes
\chapters@set{20261125T014700}{30}% pseudocharacters-and-the-length-bound
\chapters@set{20261125T011000}{31}% galois-representations-black-box
\chapters@set{20261126T030300}{32}% the-data
\chapters@set{20261126T020210}{33}% bessel-poincare-series-and-their-gram-operators
\chapters@set{20261126T030400}{34}% the-seesaw-identity
\chapters@set{20261126T030500}{35}% divisibility-of-the-composite-lift
\chapters@set{20261126T030600}{36}% theorem-a-short-forms
\chapters@set{20261126T030700}{37}% the-determinant-of-a-chain
\chapters@set{20261126T030800}{38}% the-cost-of-one-lift
\chapters@set{20261126T030900}{39}% denominators-of-projections-to-lifts
\chapters@set{20261126T031000}{40}% twists-with-few-congruences
\chapters@set{20261126T031100}{41}% theorem-c-and-the-main-inequality
\providecommand\chapnumber[1]{%
  \@ifundefined{chapters@#1}%
    {\PackageError{chapters}{Sheet #1 is not listed in the index}%
      {List it in 20261125T012000--grh-workbook__guide.org\MessageBreak
       and run scripts/make-chapters.py.}}%
    {\edef\thechap{\@nameuse{chapters@#1}}}}
\makeatother
\chapnumber{<sheet ID>}
\makeatletter
\providecommand\chapters@set[2]{\@namedef{chapters@#1}{#2}}
\chapters@set{20261125T001000}{1}% zeros-to-ramp-sums
\chapters@set{20261125T001100}{2}% inert-primes-and-estimate-b
\chapters@set{20261125T005000}{3}% galois-fourier-and-averaging
\chapters@set{20261125T005500}{4}% columns-from-a-sphere
\chapters@set{20261125T000000}{5}% slope-method
\chapters@set{20261126T030100}{6}% a-twist-for-each-prime
\chapters@set{20261125T003000}{7}% the-chain-and-the-square-root
\chapters@set{20261125T010000}{8}% rank-lemma-and-zero-estimate
\chapters@set{20261125T010500}{9}% poincare-series-and-near-orthogonality
\chapters@set{20261124T223000}{10}% elliptic-curves-for-the-proof
\chapters@set{20261125T009500}{11}% fock-model-and-theta-intertwiners
\chapters@set{20261124T223100}{12}% complex-multiplication-and-inert-reduction
\chapters@set{20261125T008000}{13}% one-cm-determinant
\chapters@set{20261124T214500}{14}% supersingular-theta-divisibility
\chapters@set{20261124T214600}{15}% the-isotropic-plane
\chapters@set{20261125T013700}{16}% pairings-on-a-formal-group-of-height-two
\chapters@set{20261125T013500}{17}% hecke-nondegeneracy-on-the-tree
\chapters@set{20261126T030200}{18}% the-bound-is-attained
\chapters@set{20261125T009000}{19}% kronecker-limit-and-chowla-selberg
\chapters@set{20261125T002000}{20}% definite-theta-lifts
\chapters@set{20261125T015500}{21}% siegel-weil-in-the-smallest-case
\chapters@set{20261125T002100}{22}% definite-theta-lifts-and-norms
\chapters@set{20261125T016500}{23}% hecke-operators-and-theta-lifts
\chapters@set{20261125T016700}{24}% divisibility-along-the-chain
\chapters@set{20261125T004000}{25}% denominators-and-sha
\chapters@set{20261125T004100}{26}% tate-twists-and-sha
\chapters@set{20261126T020430}{27}% the-minus-part-by-gauss-sums
\chapters@set{20261125T014100}{28}% congruence-modules
\chapters@set{20261125T014500}{29}% eisenstein-congruences-and-selmer-classes
\chapters@set{20261125T014700}{30}% pseudocharacters-and-the-length-bound
\chapters@set{20261125T011000}{31}% galois-representations-black-box
\chapters@set{20261126T030300}{32}% the-data
\chapters@set{20261126T020210}{33}% bessel-poincare-series-and-their-gram-operators
\chapters@set{20261126T030400}{34}% the-seesaw-identity
\chapters@set{20261126T030500}{35}% divisibility-of-the-composite-lift
\chapters@set{20261126T030600}{36}% theorem-a-short-forms
\chapters@set{20261126T030700}{37}% the-determinant-of-a-chain
\chapters@set{20261126T030800}{38}% the-cost-of-one-lift
\chapters@set{20261126T030900}{39}% denominators-of-projections-to-lifts
\chapters@set{20261126T031000}{40}% twists-with-few-congruences
\chapters@set{20261126T031100}{41}% theorem-c-and-the-main-inequality
\providecommand\chapnumber[1]{%
  \@ifundefined{chapters@#1}%
    {\PackageError{chapters}{Sheet #1 is not listed in the index}%
      {List it in 20261125T012000--grh-workbook__guide.org\MessageBreak
       and run scripts/make-chapters.py.}}%
    {\edef\thechap{\@nameuse{chapters@#1}}}}
\makeatother
\chapnumber{<sheet ID>}
\makeatletter
\providecommand\chapters@set[2]{\@namedef{chapters@#1}{#2}}
\chapters@set{20261125T001000}{1}% zeros-to-ramp-sums
\chapters@set{20261125T001100}{2}% inert-primes-and-estimate-b
\chapters@set{20261125T005000}{3}% galois-fourier-and-averaging
\chapters@set{20261125T005500}{4}% columns-from-a-sphere
\chapters@set{20261125T000000}{5}% slope-method
\chapters@set{20261126T030100}{6}% a-twist-for-each-prime
\chapters@set{20261125T003000}{7}% the-chain-and-the-square-root
\chapters@set{20261125T010000}{8}% rank-lemma-and-zero-estimate
\chapters@set{20261125T010500}{9}% poincare-series-and-near-orthogonality
\chapters@set{20261124T223000}{10}% elliptic-curves-for-the-proof
\chapters@set{20261125T009500}{11}% fock-model-and-theta-intertwiners
\chapters@set{20261124T223100}{12}% complex-multiplication-and-inert-reduction
\chapters@set{20261125T008000}{13}% one-cm-determinant
\chapters@set{20261124T214500}{14}% supersingular-theta-divisibility
\chapters@set{20261124T214600}{15}% the-isotropic-plane
\chapters@set{20261125T013700}{16}% pairings-on-a-formal-group-of-height-two
\chapters@set{20261125T013500}{17}% hecke-nondegeneracy-on-the-tree
\chapters@set{20261126T030200}{18}% the-bound-is-attained
\chapters@set{20261125T009000}{19}% kronecker-limit-and-chowla-selberg
\chapters@set{20261125T002000}{20}% definite-theta-lifts
\chapters@set{20261125T015500}{21}% siegel-weil-in-the-smallest-case
\chapters@set{20261125T002100}{22}% definite-theta-lifts-and-norms
\chapters@set{20261125T016500}{23}% hecke-operators-and-theta-lifts
\chapters@set{20261125T016700}{24}% divisibility-along-the-chain
\chapters@set{20261125T004000}{25}% denominators-and-sha
\chapters@set{20261125T004100}{26}% tate-twists-and-sha
\chapters@set{20261126T020430}{27}% the-minus-part-by-gauss-sums
\chapters@set{20261125T014100}{28}% congruence-modules
\chapters@set{20261125T014500}{29}% eisenstein-congruences-and-selmer-classes
\chapters@set{20261125T014700}{30}% pseudocharacters-and-the-length-bound
\chapters@set{20261125T011000}{31}% galois-representations-black-box
\chapters@set{20261126T030300}{32}% the-data
\chapters@set{20261126T020210}{33}% bessel-poincare-series-and-their-gram-operators
\chapters@set{20261126T030400}{34}% the-seesaw-identity
\chapters@set{20261126T030500}{35}% divisibility-of-the-composite-lift
\chapters@set{20261126T030600}{36}% theorem-a-short-forms
\chapters@set{20261126T030700}{37}% the-determinant-of-a-chain
\chapters@set{20261126T030800}{38}% the-cost-of-one-lift
\chapters@set{20261126T030900}{39}% denominators-of-projections-to-lifts
\chapters@set{20261126T031000}{40}% twists-with-few-congruences
\chapters@set{20261126T031100}{41}% theorem-c-and-the-main-inequality
\providecommand\chapnumber[1]{%
  \@ifundefined{chapters@#1}%
    {\PackageError{chapters}{Sheet #1 is not listed in the index}%
      {List it in 20261125T012000--grh-workbook__guide.org\MessageBreak
       and run scripts/make-chapters.py.}}%
    {\edef\thechap{\@nameuse{chapters@#1}}}}
\makeatother
\chapnumber{20261126T030400}
\renewcommand{\thesection}{\thechap.\arabic{section}}

\theoremstyle{definition}
\newtheorem{problem}{Problem}[section]
\theoremstyle{plain}
\newtheorem{theorem}{Theorem}[section]

\newcommand{\Z}{\mathbb{Z}}
\newcommand{\Q}{\mathbb{Q}}
\newcommand{\C}{\mathbb{C}}
\newcommand{\ii}{\mathrm{i}}
\newcommand{\Oc}{\mathcal{O}}
\newcommand{\Ac}{\mathcal{A}}
\newcommand{\Pc}{\mathcal{P}}
\newcommand{\Tc}{\mathcal{T}}
\newcommand{\Sym}{\operatorname{Sym}}
\newcommand{\Hom}{\operatorname{Hom}}
\newcommand{\Stab}{\operatorname{Stab}}
\newcommand{\ram}{\mathrm{ram}}
\newcommand{\tl}[1]{\texttt{#1}}

\setlength{\emergencystretch}{3em}

\title{The seesaw identity}
\author{}
\date{November 26, 2026}

\begin{document}
\maketitle
\refstepcounter{chap}\label{chap}

\paragraph{Prerequisites.}
From Chapter~\ref{dt-chap}: the lattices \(I_n\) over \(\Oc = \Z[\ii]\), the groups \(\Gamma_L\), forms of weight \((a)\) and their norm (Sections~\ref{dt-sec:lattices} and~\ref{dt-sec:forms}), and Problems~\ref{dt-pr:aut}, \ref{dt-pr:class-one} and~\ref{dt-pr:sym-norm}.  From Chapter~\ref{tn-chap}: the kernels \(\vartheta_\lambda(L,L')\), the lifts \(\Theta_\lambda\), the moments \(c_a\), and Problems~\ref{tn-pr:equiv}, \ref{tn-pr:adjoint}, \ref{tn-pr:ca}, \ref{tn-pr:Fa}, \ref{tn-pr:U1U1}, \ref{tn-pr:U1U2}, \ref{tn-pr:U2U1}, \ref{tn-pr:seesaw} and~\ref{tn-pr:compose}.  From Chapter~\ref{da-chap}: the admissible data, the ranks \(n_j\) and the spaces \(Y_j\).  From Chapter~\ref{bp-chap}: Bessel periods, the Poincar\'e map, and Theorem~B.  Inner products of forms on definite lattices are linear in the first variable.

\medskip
In the rewritten Lean proof \cite{GL} of Alp\"oge's theorem, Theorem~A concerns lifts of normalized Poincar\'e translates along the chain \(Y_0 \to \dots \to Y_m\), whose first space is indefinite, while the estimates of size and divisibility are proved along a chain \(X_0 \to \dots \to X_m\) of definite spaces, where forms of fixed weight and level have finitely many class coordinates.  How does the proof pass between the chains?  For each twist, the seesaw identity expresses the lift along the first chain of the Poincar\'e image of a ramified form \(\varphi\) as one nonzero constant times a combination of lifts of \(\varphi\) along the definite chain, restricted to \(Y_m\).  We first study an analogue for one lift between lattices of ranks one and two.


\section{A Poincar\'e map between definite lattices}\label{sec:toy}

What is a Poincar\'e series when the groups are finite?  Let \(X := \Oc e_1\) and \(Y := I_2 = X \perp \Oc e_2\), so that \(\Gamma_Y\) has order \(32\) and \(\Gamma_X = \mu_4\) (Problem~\ref{dt-pr:aut}).  By Problem~\ref{dt-pr:class-one} a form of weight \((a)\) on \(I_n\), \(n \le 2\), is a single vector of \(\Sym^a(\C^n)\) fixed by \(\Gamma_{I_n}\).  The inner product of forms is \(\langle\psi,\psi'\rangle := \lvert\Gamma_{I_n}\rvert^{-1}\langle\psi(I_n),\psi'(I_n)\rangle_{\Sym}\), where \(\langle\,,\rangle_{\Sym}\) is the inner product of \(\Sym^a(\C^n)\); we identify a form \(\psi\) with its value \(\psi(I_n)\).  Let \(0 \le j \le a\) be integers divisible by \(4\), and let \(\varphi_j \in \Ac_{(j)}(1)\) be the form with value \(e_1^j\).  For \(\varphi \in \Ac_{(j)}(1)\) put \(\iota\varphi := \varphi\, e_2^{a-j} \in \Sym^a(\C^2)\) and
\[
  \Pc\varphi := \sum_{\gamma \in \Gamma_Y / \Gamma_{e_2}} \gamma\, \iota\varphi ,
\]
where \(\Gamma_{e_2}\) is the stabilizer of \(e_2\) in \(\Gamma_Y\).  The restriction of \(\psi \in \Ac_{(a)}(2)\) to \(X\), paired with \(\varphi\), is the period \(B(\psi,\varphi) := \lvert\Gamma_X\rvert^{-1}\langle\psi, \iota\varphi\rangle_{\Sym}\) of Chapter~\ref{tn-chap}, Section~\ref{tn-sec:seesaw}.

\begin{problem}\label{pr:stab}
\(\Gamma_{e_2} = \{\operatorname{diag}(u,1) : u \in \mu_4\}\), and restriction to \(X\) is an isomorphism \(\Gamma_{e_2} \to \Gamma_X\).  The quotient \(\Gamma_Y/\Gamma_{e_2}\) has \(8\) elements.
\end{problem}

\begin{problem}\label{pr:poincare}
\(\Pc\varphi\) is well defined and lies in \(\Ac_{(a)}(2)\), and \(\langle \psi, \Pc\varphi\rangle = B(\psi,\varphi)\) for every \(\psi \in \Ac_{(a)}(2)\).
\end{problem}

\begin{problem}\label{pr:Pphi}
\(\Pc\varphi_j = 4\bigl(e_1^j e_2^{a-j} + e_1^{a-j} e_2^{j}\bigr)\).  For \(a = 8\): \(\Pc\varphi_4 = 8\, e_1^4 e_2^4\) and \(\Pc\varphi_0 = 4(e_1^8 + e_2^8)\).
\end{problem}


\section{One lift}\label{sec:step}

Now lift \(\Pc\varphi\) from \(Y\) to a lattice \(L' = I_{n'}\), \(n' \in \{1,2\}\), with standard basis \(e'_1, \dots, e'_{n'}\); write \(\Ac_{(d)}(L') := \Ac_{(d)}(n')\).  For \(x \colon Y \to L'\) write \(x_k := x e_k\), and also view \(x_1\) as the map \(\C e_1 \to L'_\C\), \(e_1 \mapsto x_1\).  Put
\[
  \theta_{L'} := \sum_{v \in L'} v^{a-j}\, e^{-\pi h(v,v)} \in \Sym^{a-j}(L'_\C) ,
\]
the vector obtained from the line \(\Oc e_2\) alone.  Products are taken in the symmetric algebra, as in Chapter~\ref{dt-chap}.

\begin{problem}\label{pr:split}
For \(x = (x_1, x_2)\) and \(\varphi \in \Sym^j(\C e_1)\), \(\Sym^a(x)(\varphi\, e_2^{a-j}) = \bigl(\Sym^j(x_1)\varphi\bigr)\, x_2^{a-j}\).
\end{problem}

\begin{problem}\label{pr:one-step}
For \(\varphi \in \Ac_{(j)}(1)\), with \(\Theta_{(a)} \colon \Ac_{(a)}(2) \to \Ac_{(a)}(L')\) and \(\Theta_{(j)} \colon \Ac_{(j)}(1) \to \Ac_{(j)}(L')\) the lifts of Chapter~\ref{tn-chap},
\[
  \Theta_{(a)}(\Pc\varphi) = \bigl(\Theta_{(j)}\varphi\bigr)\, \theta_{L'} .
\]
\end{problem}

The right side lifts \(\varphi\) from the smaller space \(X\), with the kernel multiplied by the fixed vector \(\theta_{L'}\).

\begin{problem}\label{pr:theta-vector}
\(\theta_{L'}\) is fixed by \(\Gamma_{L'}\).  For \(L' = I_1\), \(\theta_{I_1} = c_{a-j}\, {e'_1}^{a-j}\) and \(\Theta_{(a)}(\Pc\varphi_j) = \tfrac14 c_j c_{a-j}\, {e'_1}^a\), in agreement with Problem~\ref{tn-pr:U2U1}.  For \(L' = I_2\), \(a = 8\) and \(j = 4\),
\[
  \Theta_{(8)}(\Pc\varphi_4) = \tfrac14 c_0^2 c_4^2\,\bigl({e'_1}^8 + 2 {e'_1}^4 {e'_2}^4 + {e'_2}^8\bigr) = 4\vartheta_0^4 (\varpi/2\pi)^8\,\bigl({e'_1}^8 + 2 {e'_1}^4 {e'_2}^4 + {e'_2}^8\bigr) .
\]
\end{problem}

\begin{problem}\label{pr:transpose}
For \(\psi \in \Ac_{(a)}(L')\), let \(\Theta'_{(a)}\psi \in \Ac_{(a)}(2)\) be its lift back to \(Y\).  Then
\[
  B\bigl(\Theta'_{(a)}\psi, \varphi\bigr) = \bigl\langle \psi,\ (\Theta_{(j)}\varphi)\, \theta_{L'} \bigr\rangle .
\]
For \(L' = I_1\) this is Problem~\ref{tn-pr:seesaw}.
\end{problem}


\section{Chains, weighted sums, and translates of the vacuum}\label{sec:chain}

\begin{problem}\label{pr:chain-toy}
Let \(\Theta^{(2)}\) be the composite \(\Ac_{(8)}(2) \to \Ac_{(8)}(1) \to \Ac_{(8)}(2)\) of two lifts.  Then \(\Theta^{(2)}\Pc\varphi_4 = \Theta_{(8)}\bigl((\Theta_{(4)}\varphi_4)\,\theta_{I_1}\bigr) = 16\vartheta_0^4 (\varpi/2\pi)^{16}\bigl(17 {e'_1}^8 + 70 {e'_1}^4 {e'_2}^4 + 17 {e'_2}^8\bigr)\).
\end{problem}

By Problem~\ref{tn-pr:compose}, the kernel of \(\Theta^{(2)}\) is the sum over the middle classes of the products of the two weight-eight kernels, each divided by the order of the middle stabilizer.  In general the data of a chain of lifts form a tensor, with one factor for each lift, and the seesaw produces a particular tensor.

The general proof uses adelic unitary groups.  It pairs both sides with a form on \(Y_m\) and transfers the Bessel period along the lower chain one lift at a time.  The finite variables in the resulting outer integrals are summed over finite sets of cosets of a level, with reciprocal multiplicities as weights.  Here is a finite model for those weights.

\begin{problem}\label{pr:cut}
Let a finite group \(R\) act on a finite set \(\Omega\), let \(T \subset \Omega\) meet every orbit, and put \(\mu_T(y) := \#\{r \in R : r^{-1}y \in T\}\).  Then for every function \(f \colon \Omega \to \C\),
\[
  \sum_{y \in \Omega} f(y) = \sum_{r \in R}\, \sum_{x \in T} \frac{f(rx)}{\mu_T(rx)} .
\]
If \(f\) is \(R\)-invariant, the right side is \(\lvert R\rvert \sum_{x \in T} f(x)/\mu_T(x)\); if \(T\) has exactly one point in each orbit, then \(\mu_T(x) = \lvert\Stab_R(x)\rvert\).  For \(R = \{1,2,4,3\} \subset (\Z/5)^\times\) acting on \(\Omega = \Z/5\) by multiplication, \(T = \{0,1,2\}\) and \(f = 1\), the multiplicities are \(4, 2, 2\) and both sides equal \(5\).
\end{problem}

In the rewritten Lean proof these sums are formed for a tame character, which has the data of a twist (Section~\ref{sec:theorem}) except that \(\psi_w\) may be trivial also for \(w \in Q\); put \(q := \prod_{w \in Q}\#k_w\).  At stage \(0 \le j < m\) the sum runs over a finite set \(T_j\) of cosets \(xJ_j\), where \(J_j\) consists of the elements of the lattice stabilizer of \(Y_j\) that reduce to the identity at the places dividing \(\pi_M\) (Section~\ref{sec:theorem}) and at those in \(Q\); for \(j \ge 1\) the multiplicity of \(xJ_j \in T_j\) counts the classes of \(r\) in the finite adelic Bessel subgroup, modulo its intersection with the lattice stabilizer, with \(r^{-1}xJ_j \in T_j\), and for \(j = 0\) it is \(1\).  Each multiplicity divides some power of \(q\).

The data for the definite chain in the seesaw are finite sums of tensors built from translates of the integral vacuum of each lift.  In the simplest local model the vacuum is the characteristic function \(\mathbf 1_{M_0}\) of a lattice \(M_0 \subset \Q_p^n\), and the elementary translation operators are the shifts \(f(v) \mapsto f(v-x)\) and the multiplications \(f(v) \mapsto \chi(v)f(v)\) by additive characters \(\chi\) of \(\Q_p^n\).

\begin{problem}\label{pr:lattice}
Let \(M \subset M_0 \subset \Q_p^n\) be lattices with \([M_0 : M] = p^e\).  Then \(\mathbf 1_{M_0} = \sum_{x \in M_0/M} \mathbf 1_{x + M}\), and \(\mathbf 1_M = p^{-e} \sum_{\chi} \chi\, \mathbf 1_{M_0}\), where \(\chi\) runs over \(p^e\) characters of \(\Q_p^n\) whose restrictions to \(M_0\) are the characters of \(M_0/M\).  Hence the characteristic function of every lattice commensurable with \(M_0\) is a combination of shifts of the products \(\chi\mathbf 1_{M_0}\) with coefficients in \(\Z[1/p]\).
\end{problem}


\section{The seesaw identity}\label{sec:theorem}

We use the CM field \(E\) with \(f\) infinite places and the distinct level primes \(p_0, p_1\) of Chapter~\ref{da-chap}.  The principal level is \(M = p_0p_1 = \pi_M\bar\pi_M\) with \(\bar\pi_M = \alpha\) for the element \(\alpha \in \Oc_K\) of Chapter~\ref{da-chap}, and the level \(N = \pi_N\bar\pi_N\), \(\pi_N \in \Oc_E\), is prime.  Write \(\varrho := m/2 - 1\), where \(m\) is the number of lifts.  The lower chain starts at the indefinite space \(Y_0\) and ends at the definite space \(Y_m\).  Let \(X_j\) be the standard definite hermitian space over \(E\) of rank \(n_j\), for \(j \ge 0\); the last grading identifies \(Y_m\) with the span of the first \(n_m - 1\) basis vectors of \(X_m\).  The following objects are those of the rewritten Lean proof \cite{GL}.

\emph{Twists.}  A twist \(\eta\) consists of an automorphic character of the norm-one torus of \(E/E^+\), a finite set \(Q\) of finite places of \(E\), and a multiplicative character \(\psi_w\) of the residue field \(k_w\) at each finite place \(w\).  The character has weight zero, and its value at every norm-one idele \(u\) with integral finite part is \(\prod_w \psi_w(u_w \bmod w)\), with \(\psi_w\) trivial for \(w \notin Q\) and nontrivial for \(w \in Q\); moreover \(Q\) contains the conjugate of none of its places, and \(q := \prod_{w \in Q} \#k_w\) is prime to \(MN\).

\emph{Source forms.}  The source forms \(S\) are the algebraic forms on \(X_0\) fixed by its level.  At an infinite place \(\tau\) the Weyl module of their values has highest weight \((a_\tau, 0, \dots, 0)\), where \(a_\tau = a = xb\) if \(\tau \equiv 1 \bmod q_0\) and \(a_\tau = 0\) otherwise.  The level consists of the finite adelic unitary matrices preserving the standard lattice that reduce to the identity at each \(w \mid \pi_M\) and whose last row is congruent to \((0, \dots, 0, 1)\) modulo \(\pi_N\).  For each place \(w \mid \pi_N\) a finite idele \(z_w\) of norm one is fixed: it is integral, equal to \(1\) outside \(w\) and \(\bar w\), and of order \(N-1\) modulo \(w\).  It acts on \(S\) by right translation by the scalar \(z_w\), an operator \(T_w\) with \(T_w^{N-1} = 1\).  A source form \(\varphi\) is \emph{ramified} if \(\sum_{i=0}^{N-2} T_w^i\varphi = 0\) for every \(w \mid \pi_N\); these form a subspace \(S_\ram\).

\emph{The lower chain.}  For a point \(g\) of the archimedean unitary group of \(Y_0\) and a linear form \(t\) on the values of source forms, the Poincar\'e map \(\Pc_{g,t}\) from \(S\) to the forms on \(Y_0\) of the principal level and the discrete series type of Chapter~\ref{bp-chap} is defined by \(\langle y, \Pc_{g,t}\varphi\rangle = \int_{[R_0]} W_{t,\varphi}(r)\,\overline{y(r(g,1))}\,dr\) for all such \(y\), where the Petersson product is conjugate linear in the first variable, \(R_0\) is the Bessel subgroup of \(Y_0\), and \(W_{t,\varphi}\) is the Bessel function of \(\varphi\) formed with \(t\).  For \(1 \le i \le m\) let \(\mu_i^0\) be the fixed untwisted character of the step \(Y_{i-1} \to Y_i\), and \(\upsilon\) the fixed splitting character restricted to the norm-one torus.  For a tame character \(\eta\), the lift of this step uses on each of its two spaces, of rank \(d\), the character \(\mu_i^0\eta^i\upsilon^{\lfloor d/2\rfloor}\), its datum is the prescribed joint harmonic at infinity times the finite point datum given by the inclusion matrix at the level places and the vacuum elsewhere, and it is multiplied by the index of \(J_{i-1}\) in the source lattice stabilizer.  The composite of the \(m\) lifts, followed by the passage to algebraic forms on \(Y_m\), is \(L_\eta\).

\emph{The definite chain.}  Let \(\Tc\) be the set of tuples \(x = (x_i)_{i \ge 0}\) of matrices \(x_i\) of finite adeles of \(E\), with rows indexed by the basis of \(X_{i+1}\) and columns by that of \(X_i\).  The datum \(v_i(x)\) is the integral vacuum of \(\Hom(X_i, X_{i+1})\) translated by the finite Heisenberg operator of \(x_i\).  For a finitely supported \(\Psi \colon \Tc \to \C\), the glued tensor is
\[
  \Psi^{\mathrm{gl}} := \sum_{x} \Psi(x) \bigotimes_{i=0}^{\varrho} G_i\bigl(v_{2i}(x) \otimes v_{2i+1}(x)\bigr) ,
\]
where \(G_i\), for the double lift \(X_{2i} \to X_{2i+1} \to X_{2i+2}\), sums over the vectors \(y\) of residues modulo \(N\), each represented in \(\{0, \dots, N-1\}\), indexed by a basis of \(X_{2i+1}\) and the places \(w \mid \pi_N\): it translates the first factor by the last column \(-y/N\) supported at these places, and the second by the conjugate last row supported at their conjugates.  The composite lift \(B_{\eta,\Psi}(\varphi)\) applies the \(\varrho + 1\) double lifts with data \(\Psi^{\mathrm{gl}}\) and twist \(\eta\) to \(\varphi\), and restricts the result from \(X_m\) to \(Y_m\), by evaluating at \(\operatorname{diag}(h,1)\), applying a fixed branching map, and multiplying by \(\prod_\tau \Phi_\tau(\bar\pi_N)^{m r^8 - a_\tau}\), where \(\Phi_\tau\) is the embedding of the CM type at \(\tau\) and \(r\) is the rank of \(Y_0\).

\emph{Chain data.}  \(\Psi\) is a chain datum for the twist \(\eta\) if: (i) for \(x\) in the support of \(\Psi\), every entry of every \(x_i\) becomes integral at every finite place \(w \nmid q\) after multiplication by \(\bar\pi_M\), and vanishes at every \(w \nmid Mq\); (ii) every value \(\Psi(x)\) is integral away from \(q\): \(\lvert\iota(\Psi(x))\rvert_\ell \le 1\) for every prime \(\ell \nmid q\) and every ring homomorphism \(\iota \colon \C \to \C_\ell\); (iii) if \(Q = \emptyset\), then \(\sum_x \lvert\sigma(\Psi(x))\rvert \le M^{3fmn_m^2}\) for every field automorphism \(\sigma\) of \(\C\); (iv) \(\Psi^{\mathrm{gl}}\) is fixed by the levels of \(X_0, \dots, X_{m-1}\), acting through the Weil representations twisted by \(\eta\), with the identity on \(X_m\).

\begin{theorem}[the seesaw identity]\label{thm:seesaw}
Let \(D\) be admissible data, \(t\) a linear form on the values of the source forms, and \(g\) a point of the archimedean unitary group of \(Y_0\).  There is \(c \in \C^\times\) such that for every twist \(\eta\) there is a finitely supported \(\Psi \colon \Tc \to \C\) which is a chain datum for \(\eta\) and satisfies
\[
  L_\eta\bigl(\Pc_{g,t}\varphi\bigr) = c\, B_{\eta,\Psi}(\varphi) \qquad\text{for every } \varphi \in S_\ram .
\]
\end{theorem}

For fixed \(D\), \(t\), \(g\), one scalar \(c\) works for all twists, and for each twist one \(\Psi\) works for all \(\varphi \in S_\ram\); the equality is one of algebraic forms on \(Y_m\).

\paragraph{In the Lean formalization.}
{\small\raggedright Theorem~\ref{thm:seesaw} is \tl{New.Data.seesaw\_identity} in \tl{New/UpperBound/Seesaw/SeesawIdentity.lean} of the rewritten Lean proof \cite{GL}, with \tl{IsChainDatum}, \tl{gluedTensor} and \tl{compositeLift} in \tl{New/UpperBound/Objects/Columns.lean}.  The steps of its proof described below are \tl{integral\_mul\_liftY}, \tl{exists\_integral\_mul\_lowerChainUpTo\_eq\_cutSum}, \tl{lowerChain\_eq\_of\_lowerChainUpTo\_eq}, \tl{exists\_lowerChainUpTo\_poincareAt\_eq} and \tl{exists\_isAdmissibleCombination\_seesawDatum} in \tl{SeesawIdentity.lean}, with the divisibility of the multiplicities given by \tl{exists\_cutMultiplicity\_dvd} in the adjacent \tl{Cut.lean}; the statements on ramified forms in Section~\ref{sec:claim} are \tl{ramifiedSourceForms\_eq\_range}, \tl{tensorChainLift\_eq\_zero\_of\_ramified}, \tl{doubleLiftChain\_eq\_zero\_of\_ramified}, \tl{unfoldStage\_eq\_zero\_of\_scalingGenerator} and \tl{hasLatticeBounds\_of\_unfoldStage\_ne\_zero}, and the application in Section~\ref{sec:claim} is \tl{exists\_lowerChain\_poincare\_eq}.  The other objects of this chapter are, in order: Theorem~A, \tl{span\_short\_integralRamified\_eq\_top} and \tl{indexLe\_normalizedPoincare}, with the definite size and divisibility estimates \tl{weylNormSq\_compositeLift\_le} and \tl{exists\_err\_gaugeLe\_compositeLift}, and finiteness of the class sets from \tl{hasFiniteType\_of\_isLevel}; the tensor \tl{seesawDatum} in \tl{ChainTensor}, and \tl{IsChainDatum.exists\_isWeakChainDatum\_compositeLift\_eq}; the sums \tl{cutSumAt}, \tl{cutSum}, \tl{cutPoints}, \tl{cutMultiplicity}, \tl{coverMultiplicity}, \tl{levelY} and \tl{IsStageCut.ne\_zero\_of\_hasBoundedOrbit}, with the unfolding step \tl{integral\_unfoldStageOn\_mul\_transposedLiftY}; the spaces \tl{homX}, \tl{spaceY\_gram\_eq}, \tl{lastGrading}, \tl{m\_eq\_two\_mul\_varrho\_add\_two}; the characters \tl{TameChar}, \tl{TameChar.normQ}, \tl{Twist}, \tl{AutChar.HasWeight}; the source \tl{SourceForm}, \tl{weightX\_eq}, \tl{levelX}, \tl{scalingGenerator}, \tl{centralRep}, \tl{centralRep\_scalingGenerator\_pow\_eq\_one}, \tl{ramifiedSourceForms}; the lower chain \tl{poincareAt}, \tl{inner\_poincareAt}, \tl{poincareAt\_smul}, \tl{lowerChain}, \tl{liftY}, \tl{datumY}, \tl{thetaKernelY}; the upper data \tl{translateDatum}, \tl{Translations}, \tl{intVacuum}, \tl{doubleLiftGluing}, \tl{restrictToY}, \tl{branchExponent\_eq}, \tl{IsAdmissibleCoeffs}, \tl{integralAway}, \tl{ChainTensor.IsInvariant}, \tl{residueMatrixCount}, \tl{IsAdmissibleCombination.normUnit\_smul}; and for the application \tl{poincare}, \tl{seesawConst}, \tl{normalizedPoincare}, Theorem~B \tl{enough\_forms}, and \(b > 2\) from \tl{le\_constantBound} and the scale inequalities of \tl{New.Data}.\par}

\medskip
The rewriting proves Theorem~\ref{thm:seesaw} along the lines of Section~\ref{sec:step}: it pairs both sides with a form on \(Y_m\), transfers the Bessel period along the lower chain one lift at a time by transposition as in Problem~\ref{pr:transpose} and unfolding, and sums the finite variables over the cosets described after Problem~\ref{pr:cut}.  This gives \(c_0 \in \C\), possibly zero, such that for every tame character \(\eta\) and every \(\varphi \in S_\ram\), the composite of the \(m\) lower lifts of \(\Pc_{g,t}\varphi\), as an automorphic form, is \(c_0\) times the restriction to \(Y_m\) of the chain of lifts of \(\varphi\) with a certain tensor \(\varphi_\eta\) of data; passing to algebraic forms gives the corresponding equality for \(L_\eta\).  There is \(C \in \C\), independent of \(\eta\), such that \(\varphi_\eta\) is \(C\) times a finite combination of the glued vacua \(\bigotimes_i G_i(v_{2i}(x) \otimes v_{2i+1}(x))\) whose coefficients, among other conditions, become integral away from \(q\) after multiplication by \(M^{2mfn_m^2}\).  The constant \(c\) of the theorem is then produced as in Problem~\ref{pr:unit} below, applied to \(c_0C/M^{2mfn_m^2}\).


\section{Ramified forms, and what is not claimed}\label{sec:claim}

Where does the proof use the hypothesis that \(\varphi\) is ramified?

\begin{problem}\label{pr:mean}
Let \(n \ge 1\) be an integer and let \(T\) be a linear endomorphism of a finite-dimensional complex vector space \(V\) with \(T^n = 1\).  Then \(\ker\bigl(\sum_{i=0}^{n-1} T^i\bigr) = \operatorname{im}(1 - T)\), the sum of the eigenspaces of \(T\) for eigenvalues \(\ne 1\).  For the cyclic shift on \(\C^n\) this space has dimension \(n - 1\).
\end{problem}

\begin{problem}\label{pr:mean-several}
Let \(n \ge 1\) be an integer, and let \(T_1, \dots, T_s\) be commuting linear endomorphisms of a finite-dimensional complex vector space \(V\) with \(T_k^n = 1\) for every \(k\).  Then \(\bigcap_k \ker\bigl(\sum_{i=0}^{n-1} T_k^i\bigr) = \operatorname{im}\prod_k (1 - T_k)\).
\end{problem}

So \(S_\ram\) is the image of \(\prod_{w \mid \pi_N}(1 - T_w)\).

\begin{problem}\label{pr:fixed}
In the situation of Problem~\ref{pr:mean-several}, let \(\Lambda \colon V \to V'\) be linear with \(\Lambda T_k = \Lambda\) for one \(k\).  Then \(\Lambda\) vanishes on \(\bigcap_k \ker\bigl(\sum_{i=0}^{n-1} T_k^i\bigr)\).
\end{problem}

The rewriting uses Problem~\ref{pr:fixed} for chains of lifts: if the Weil action of \((1, z_w)\) on one lift factor alone fixes the tensor of data of a chain, the chain annihilates \(S_\ram\).  In the unfolding, at the even stages \(2 \le j \le m-2\), the function left by unfolding a ramified form vanishes at the points \((a, x)\), and their translates by the Bessel subgroup, with \(x^{-1} z_w x\) in the level of \(Y_j\), where \(z_w\) acts as a scalar on \(X_j\) and trivially on its complement.  If the initial finite point lies in the principal level and the unfolded function at stage \(1 \le j < m\) is nonzero at \((a,x)\), an induction based on this vanishing bounds, for every finite adelic point \(r'\) of the Bessel subgroup for which the lattice of \(r'x\) is graded, the intersection of that lattice with the complement of \(X_j\); so the orbit of the level coset of \(x\) meets the finite set of cosets used at stage \(j\).

Theorem~\ref{thm:seesaw} asserts that \(c \ne 0\), but not that either side is nonzero, for any \(\varphi\) or any twist.

\begin{problem}\label{pr:unit}
Let \(V, V'\) be complex vector spaces, \(A, B \colon V \to V'\) linear maps, and \(c_0 \in \C\) such that \(A = c_0 B\).  Put \(u := c_0^{-1}\) if \(c_0 \ne 0\) and \(u := 1\) otherwise.  Then \(u \ne 0\) and \(A = u^{-1}(u c_0 B)\), and \(uc_0B = 0\) if and only if \(c_0 = 0\) or \(B = 0\).
\end{problem}

\begin{problem}\label{pr:degenerate}
For \(t = 0\), Theorem~\ref{thm:seesaw} holds with \(c = 1\) and \(\Psi = 0\) for every twist: \(\Pc_{g,0} = 0\), the function \(\Psi = 0\) is a chain datum, and \(B_{\eta,0} = 0\).
\end{problem}

The proof of Theorem~A applies Theorem~\ref{thm:seesaw} to the Poincar\'e map \(\Pc\) of Chapter~\ref{bp-chap}: the lower chain applied to \(c^{-1}\Pc\varphi\) equals the restricted definite chain applied to \(\varphi\), for \(\varphi \in S_\ram\).  Write \(P_s\) for the normalized projected translates of \(\Pc\), restricted to \(S_\ram\), and \(d_T\) for the dimension of the good target forms.  Theorem~B of Chapter~\ref{bp-chap} gives \(d_T > 0\) and a finite set \(U\) in the level group indexing the translates such that \(\#U \dim S_\ram \le 5 d_T \log(1 + 2b^2)\) and \(\dim \sum_{s \in U} P_s(S_\ram) \ge (1 - 4/b^2)\, d_T\).  Since \(b > 2\), some \(P_s\) is nonzero, and each \(P_s\) factors through \(\Pc\), so \(\Pc\) is not zero on \(S_\ram\).  Theorem~B gives no nonvanishing after the lower chain of lifts.


\section*{Hints}

\subsection*{\S\ref{sec:toy}}
\begin{description}
\item[\ref{pr:poincare}] \(\gamma\) is unitary and fixes \(\psi\); count the cosets.
\item[\ref{pr:Pphi}] Use the representatives \(\operatorname{diag}(1,u)\) and \(\left(\begin{smallmatrix}0&u\\1&0\end{smallmatrix}\right)\), \(u \in \mu_4\), and \(4 \mid a - j\).
\end{description}

\subsection*{\S\ref{sec:step}}
\begin{description}
\item[\ref{pr:split}] \(e_1^j e_2^{a-j}\) is the symmetrization of \(e_1^{\otimes j} \otimes e_2^{\otimes(a-j)}\), and \(\Sym^a(x)\) is the restriction of \(x^{\otimes a}\).
\item[\ref{pr:one-step}] Problems~\ref{tn-pr:equiv} and~\ref{pr:split}; \(h(x,x) = h(x_1,x_1) + h(x_2,x_2)\) and \(\lvert\Gamma_{e_2}\rvert = \lvert\Gamma_X\rvert\).
\item[\ref{pr:theta-vector}] Substitute \(v \mapsto \gamma v\).  For the values use Problems~\ref{tn-pr:U1U1}, \ref{tn-pr:U1U2}, \ref{tn-pr:ca} and~\ref{tn-pr:Fa}, and \(c_k = 0\) for \(4 \nmid k\).
\item[\ref{pr:transpose}] Problems~\ref{pr:poincare} and~\ref{tn-pr:adjoint}.
\end{description}

\subsection*{\S\ref{sec:chain}}
\begin{description}
\item[\ref{pr:chain-toy}] Apply the second lift to both sides of Problem~\ref{pr:one-step}; then Problems~\ref{pr:theta-vector} and~\ref{tn-pr:U1U2}.
\item[\ref{pr:cut}] Exchange the two sums on the right and group by \(y = rx\).
\item[\ref{pr:lattice}] Orthogonality of the characters of the finite group \(M_0/M\).  Characters of \(M_0\) extend to \(\Q_p^n\).  For a general lattice \(M'\), use \(M_0 \cap M'\).
\end{description}

\subsection*{\S\ref{sec:claim}}
\begin{description}
\item[\ref{pr:mean}] \(T\) is diagonalizable; for Problem~\ref{pr:mean-several} use simultaneous eigenspaces.
\item[\ref{pr:degenerate}] The composite lift is linear in the glued tensor.
\end{description}

\begin{thebibliography}{GL}
\bibitem[GL]{GL} L.~Sernau, \emph{GRH for Dirichlet \(L\)-functions, in Lean~4}, a shorter rewriting of the Lean proof, text by Claude, \url{https://github.com/lsernau/grh-lean}, commit \texttt{72fead1}, readable edition \texttt{grh/}.
\end{thebibliography}

\end{document}
